\documentclass[a4paper, 11pt]{article}
\usepackage{xeCJK}
\usepackage{geometry}
\usepackage{titlesec}
\usepackage{graphicx}   % Required for images
\usepackage{subcaption}
\usepackage{paracol}
\usepackage[backend=biber, style=numeric]{biblatex}
\addbibresource{../../ref.bib}

\geometry{margin=1in}
\setCJKmainfont{MS Mincho}

\titlespacing*{\section}{0pt}{2ex}{1ex}
\titlespacing*{\subsection}{0pt}{1ex}{0.8ex}
\titlespacing*{\subsubsection}{0pt}{1ex}{0.5ex}

\begin{document}

\begin{center}
    {\LARGE \bfseries 進捗報告} \\ % Title: Large and bold
    \vspace{0.3em} % A small vertical space
    Juan Pablo Corredor \\ % Author
    ２０２６年１月２３日 % Date
\end{center}
\vspace{0.5em} % Adjust space between title block and main text

\section*{近況}
% Write your recent findings here.
\begin{itemize}
    \item まだかなり忙しい！
    \item 寒い寒い！
\end{itemize}

\section*{やったこと}
\begin{itemize} 
    \item 期末課題ｖ２。
    \item 音楽音響研究会に申し込み。
    \item 
\end{itemize}

\section*{音楽研究会の論文について}

\begin{paracol}{2}
    \subsection*{English}
    Pitch-Agnostic Partial Tracking and Low Dimensional Timbre Representation of Plucked Strings for Additive Resynthesis
    \switchcolumn

    \subsection*{日本語}
    部分音追跡による低次元音色表現：撥弦楽器のためのピッチに依存しない手法
    ピッチに依存しない部分音追跡による低次元音色表現
    ピッチに依存しない部分音追跡による低次元音色表現－撥弦楽器を例にー
    \switchcolumn*

    This paper proposes a spectral analysis and resynthesis framework for plucked string instruments, building upon the partial tracking methods described by Caetano in \cite{soundMorphingCaetano}. We introduce a decomposition system that utilizes the Continuous Wavelet Transform (CWT) for high-resolution transient analysis and Harmonic Ratio Envelopes for the sustained component of recorded samples. Unlike traditional absolute-frequency tracking, this method normalizes partial trajectories relative to the fundamental frequency ($f_0$). This yields a pitch-agnostic, low-dimensional representation that significantly compresses the spectral data with a compression ratio of approximately 4:1 (75\% reduction) while maintaining high reconstruction fidelity. This structural decoupling of pitch and timbre allows for direct statistical comparisons between different plucked string instruments and enables flexible additive resynthesis and timbre interpolation.
    \switchcolumn

    本稿では、Caetano \cite{soundMorphingCaetano} による部分音追跡法を発展させた、撥弦楽器のためのスペクトル解析および再合成手法を提案する。提案手法では、トランジェント成分の高精度な解析に連続ウェーブレット変換 (CWT) を用い、持続音成分には録音サンプルの「倍音比エンベロープ (Harmonic Ratio Envelopes)」を用いる分解システムを導入する。従来の絶対周波数に基づく追跡とは異なり、本手法は部分音の軌跡を基本周波数 ($f_0$) で正規化する。これにより、高い再合成忠実度を維持しつつ、スペクトルデータを約4:1（75\%削減）の圧縮率で大幅に削減可能な、ピッチ非依存かつ低次元な表現が得られる。このピッチと音色の構造的な分離により、異なる撥弦楽器間での直接的な統計的比較が可能となり、柔軟な加算再合成および音色補間が実現される。
    \switchcolumn*

\end{paracol}

\section*{やること}
% Outline future steps here.    
\begin{itemize}
    \item ウェーブレット変換の長さを変更して、結果を比較。
    \item モデル開発。
\end{itemize}

%開放弦

% Add your bibliography here.
\printbibliography

\end{document}